\documentclass[11pt]{article}

\usepackage[margin=1in]{geometry}
\usepackage{amsmath}
\usepackage{amssymb}
\usepackage{booktabs}
\usepackage{enumitem}
\usepackage{hyperref}

\title{Reinforcement Learning Framework for Adaptive DBS in the STN-GPe System}
\author{}
\date{}

\begin{document}
\maketitle

\section{Goal}

The goal is to train an adaptive deep brain stimulation (DBS) controller for the
STN-GPe spiking network. The controller should:

\begin{itemize}[nosep]
    \item reduce synchrony in the STN network,
    \item increase spectral entropy of the STN local field potential (LFP),
    \item use low stimulation energy.
\end{itemize}

In simple terms, the controller should push the Parkinsonian network away from a
highly synchronized, low-entropy state while avoiding unnecessary stimulation.

\section{STN-GPe System}

The simulator is a two-population spiking model with STN and GPe neurons. The
network evolves at a small integration time step:

\[
\Delta t = 0.1 \text{ ms}.
\]

The original simulator runs a full simulation in one call. For reinforcement
learning, we use a steppable wrapper. This wrapper keeps the simulator state in
memory and advances the model only for the biological time selected by the
agent.

\section{MDP Definition}

We model the control problem as a Markov decision process (MDP):

\[
\mathcal{M} = (\mathcal{S}, \mathcal{A}, P, r, \gamma).
\]

At each decision step \(t\), the agent observes a state \(s_t\), chooses an
action \(a_t\), receives a reward \(r_t\), and the simulator moves to the next
state \(s_{t+1}\):

\[
s_t \xrightarrow{a_t} s_{t+1}, \qquad r_t = r(s_t, a_t, s_{t+1}).
\]

Because the agent chooses the time until the next pulse, one RL step can span a
variable amount of biological time. This is closer to a semi-MDP, but we still
use the standard Gymnasium interface:

\[
(s_{t+1}, r_t, \text{done}, \text{info}) = \text{env.step}(a_t).
\]

\subsection{State}

The environment state is built from recent STN network metrics over a rolling
window. The observation is:

\[
s_t =
\begin{bmatrix}
R_t \\
B_t \\
H_t \\
\bar{a}_{t-1}
\end{bmatrix}.
\]

Here:

\begin{itemize}[nosep]
    \item \(R_t\) is STN synchrony. Lower is better.
    \item \(B_t\) is beta-band power. It is observed and logged, but it is not
    used in the current reward by default.
    \item \(H_t\) is spectral entropy of the STN LFP. Higher is better.
    \item \(\bar{a}_{t-1}\) is the previous action, normalized to the range
    \([0, 1]\). This helps the policy know what stimulation was recently used.
\end{itemize}

The current metric window is:

\[
T_{\text{window}} = 0.25 \text{ s}.
\]

With \(\Delta t = 0.1\) ms, this is:

\[
2500 \text{ simulator steps}.
\]

\subsection{Action}

The current action space is continuous. The agent controls:

\[
a_t =
\begin{bmatrix}
A_t \\
T_t
\end{bmatrix},
\]

where:

\begin{itemize}[nosep]
    \item \(A_t\) is the DBS pulse amplitude.
    \item \(T_t\) is the pulse period, meaning the time from this pulse to the
    next pulse.
\end{itemize}

The current amplitude range is:

\[
A_t \in [0, 200].
\]

The current frequency range is:

\[
f_t \in [5, 50] \text{ Hz}.
\]

Therefore the pulse period is bounded by:

\[
T_t = \frac{1000}{f_t} \text{ ms},
\]

so:

\[
T_t \in [20, 200] \text{ ms}.
\]

Each action emits one charge-balanced biphasic pulse:

\[
I_{\text{DBS}}(t) =
\begin{cases}
+A_t, & \text{positive phase}, \\
0, & \text{interphase gap}, \\
-A_t, & \text{negative phase}.
\end{cases}
\]

The current phase shape is fixed:

\[
\text{phase width} = 0.2 \text{ ms}, \qquad
\text{interphase gap} = 1.0 \text{ ms}.
\]

Only amplitude and pulse period are controlled by the agent at the moment.

\subsection{Transition}

After the agent chooses \(a_t\), the environment:

\begin{enumerate}[nosep]
    \item creates one biphasic DBS pulse using \(A_t\),
    \item inserts that pulse at the start of the decision window,
    \item advances the STN-GPe simulator for \(T_t\) ms,
    \item recomputes the rolling network metrics,
    \item returns the next state \(s_{t+1}\).
\end{enumerate}

Mathematically, the transition is:

\[
s_{t+1} = P(s_t, a_t),
\]

where \(P\) is the STN-GPe simulator dynamics plus the DBS input.

\section{Reward Function}

The current reward is designed to match the project objective directly:

\[
r_t = w_H H_t - w_R R_t - \lambda_q q_t - p_t.
\]

The current weights are:

\[
w_H = 1.0, \qquad w_R = 2.0, \qquad \lambda_q = 0.005.
\]

Here:

\begin{itemize}[nosep]
    \item \(H_t\) is spectral entropy, so higher entropy increases reward.
    \item \(R_t\) is synchrony, so higher synchrony decreases reward.
    \item \(q_t\) is absolute stimulation charge, so more energy decreases
    reward.
    \item \(p_t\) is an extra bad-state penalty.
\end{itemize}

The charge term is computed using absolute current:

\[
q_t = \sum_k |I_{\text{DBS}}[k]| \Delta t.
\]

Even though the biphasic pulse has net zero charge, we still penalize absolute
charge because both phases consume stimulation energy.

\subsection{Bad-State Penalty}

We add an extra penalty when the network is clearly in a bad state:

\[
p_t =
\begin{cases}
1.0, & R_t > 0.5 \text{ and } H_t < 0.5, \\
0, & \text{otherwise}.
\end{cases}
\]

Therefore the reward becomes:

\[
r_t =
H_t - 2R_t - 0.005q_t
- \mathbb{I}(R_t > 0.5 \ \text{and}\ H_t < 0.5).
\]

This means the agent receives a stronger penalty when synchrony is high and
entropy is low at the same time.

\section{Episode Structure}

Each episode has two parts.

\subsection{Warmup}

The network first runs without stimulation:

\[
T_{\text{warmup}} = 0.25 \text{ s}.
\]

This lets the network settle and fills the rolling metric buffer. Warmup steps
are not counted as RL reward steps.

\subsection{Control Period}

After warmup, the agent controls the DBS pulses:

\[
T_{\text{control}} = 1.0 \text{ s}.
\]

Because the agent can choose periods between 20 ms and 200 ms, the number of RL
decisions per episode is approximately:

\[
5 \leq N_{\text{decisions}} \leq 50.
\]

\section{Return and Best Model}

The return of an episode is the sum of rewards:

\[
G = \sum_{t=0}^{T-1} r_t.
\]

During training, evaluation is run every fixed number of environment steps. The
best model is the one with the highest mean evaluation return:

\[
\text{score} =
\frac{1}{N_{\text{eval}}}
\sum_{i=1}^{N_{\text{eval}}} G_i.
\]

The current evaluation setting uses:

\[
N_{\text{eval}} = 3.
\]

So the saved best model is not selected from synchrony alone or entropy alone.
It is selected using the combined scalar reward, which includes entropy,
synchrony, charge, and the bad-state penalty.

\section{Algorithms}

The code supports three Stable-Baselines3 algorithms.

\subsection{SAC}

Soft Actor-Critic (SAC) is the current recommended algorithm. It is useful here
because:

\begin{itemize}[nosep]
    \item the action space is continuous,
    \item the simulator is expensive,
    \item SAC is off-policy and can reuse past experience from a replay buffer,
    \item SAC encourages exploration through an entropy term.
\end{itemize}

SAC learns a policy:

\[
a_t = \pi_\theta(s_t),
\]

and updates it to maximize expected return while maintaining exploration.

\subsection{PPO}

Proximal Policy Optimization (PPO) is also supported. PPO is an on-policy
method. It is usually stable, but it can require more simulator samples because
it does not reuse old data as directly as SAC.

\subsection{TD3}

Twin Delayed Deep Deterministic Policy Gradient (TD3) is also supported for
continuous control. It uses deterministic actions plus exploration noise. It may
work, but SAC is usually a better first choice for this system because SAC has
stronger built-in exploration.

\section{Files in the Codebase}

\begin{table}[h]
\centering
\begin{tabular}{ll}
\toprule
\textbf{File} & \textbf{Purpose} \\
\midrule
\texttt{rl\_stn\_gpe/env.py} & Gymnasium environment for STN-GPe control \\
\texttt{rl\_stn\_gpe/reward.py} & Reward calculation \\
\texttt{rl\_stn\_gpe/config.py} & Action space, timing, reward, algorithm choice \\
\texttt{rl\_stn\_gpe/stepper.py} & Steppable wrapper around the simulator \\
\texttt{rl\_stn\_gpe/train.py} & Training script \\
\texttt{rl\_stn\_gpe/eval.py} & Evaluation and plotting script \\
\texttt{rl\_stn\_gpe/hparams/sac.py} & SAC hyperparameters \\
\texttt{rl\_stn\_gpe/hparams/ppo.py} & PPO hyperparameters \\
\texttt{rl\_stn\_gpe/hparams/td3.py} & TD3 hyperparameters \\
\bottomrule
\end{tabular}
\end{table}

\section{Summary}

The RL controller observes STN network metrics, chooses adaptive DBS pulse
amplitude and timing, and receives reward for lowering synchrony, increasing
entropy, and using less stimulation. SAC is the main algorithm because the
control problem is continuous and simulator samples are expensive. The best
model is saved based on the highest mean evaluation return under the current
reward function.

\end{document}
